\documentclass[aspectratio=169,10pt]{beamer}
\usepackage[UTF8]{ctex}
\usetheme{Madrid}
\usecolortheme{default}
\setbeamertemplate{navigation symbols}{}
\setbeamertemplate{footline}[frame number]
\usepackage{amsmath,amssymb,amsthm,booktabs,mathtools}
\usepackage{graphicx,tikz}
\usetikzlibrary{arrows.meta,positioning,decorations.pathreplacing,shapes,fit}

\definecolor{positive}{HTML}{0173B2}
\definecolor{negative}{HTML}{DE8F05}
\definecolor{emphasis}{HTML}{029E73}
\newcommand{\pos}[1]{\textcolor{positive}{#1}}
\newcommand{\con}[1]{\textcolor{negative}{#1}}
\newcommand{\HL}[1]{\textcolor{emphasis}{#1}}

\title[四旋翼视觉避障]{基于神经网络的四旋翼飞行器\\端到端视觉避障}
\subtitle{Mamba架构与跨架构知识蒸馏的系统性研究}
\author[邢锦文]{邢锦文 \quad 指导教师：芦安洋教授}
\institute[东北大学]{信息科学与工程学院·自动化\\东北大学}
\date{}

\begin{document}

\begin{frame}\titlepage\end{frame}

% ═══════ 1. THE SETUP ═══════
\begin{frame}{这个研究是怎么开始的}
\begin{columns}[T]
\column{0.5\textwidth}
\begin{block}{一个很朴素的问题}
\textbf{Mamba能不能替代Transformer？}\\[4pt]
\footnotesize
Transformer：O(n$^2$)，精度高但慢\\
Mamba：O(n)，理论更快，但在避障任务上\\
到底表现如何？——\textit{没有人系统测过。}
\end{block}

\begin{block}{任务场景}
四旋翼高速避障 (5--7m/s)\\
深度图 60$\times$90 $\to$ 3D速度指令\\
\footnotesize 基于ViT-Fly框架 (ICRA 2025)
\end{block}

\column{0.5\textwidth}
\begin{alertblock}{第一阶段：自动化探索}
\begin{enumerate}
\item 训练管线自动从GitHub拉取Mamba变体
\item 自动编写模型代码 $\to$ BC训练
\item \con{结果：效果不好}
\end{enumerate}
\end{alertblock}

\begin{block}{当时面临的问题}
\footnotesize
既然BC效果不好，\textbf{怎么让它好起来？}\\
既要公平对比ViT基线（同用模仿学习），\\
又要想办法提升Mamba的表现。
\end{block}
\end{columns}
\end{frame}

% ═══════ 2. THE PIVOT ═══════
  \begin{frame}{转折点：从「怎么比」到「怎么让它更好」}
  \begin{columns}[T]
  \column{0.48\textwidth}
  \begin{block}{AI提出的方向}
  \footnotesize
  既然MOHAWK证明了Transformer能蒸馏到SSM，\\
  要不要把ViT教师的知识蒸馏到你的Mamba学生？
\end{block}
\vspace{0.3cm}
\begin{block}{为什么接受这个建议}
\begin{itemize}\footnotesize
\item 保持了与ViT基线的公平性（同用模仿学习框架）
\item 蒸馏是提升性能的已知有效路径
\item 但——\HL{跨架构蒸馏在避障中从未被验证过}
\end{itemize}
\end{block}

\column{0.48\textwidth}
\begin{block}{新的研究问题}
\footnotesize
\begin{enumerate}
\item \textbf{哪个Mamba变体}在避障中表现最好？
\item ViT教师能\textbf{成功蒸馏}到Mamba学生吗？
\item 如果能，\textbf{什么时候有效，什么时候失效？}
\end{enumerate}
\end{block}
\vspace{0.3cm}
\begin{alertblock}{研究范式的转变}
\footnotesize
从「自动化管线生成模型」转向\\
「AI建议方向 + 人类决策 + 系统实验验证」
\end{alertblock}
\end{columns}
\end{frame}

% ═══════ 3. WHAT WE BUILT ═══════
\begin{frame}{我们做了什么：六种变体 + 蒸馏框架}
\begin{columns}[T]
\column{0.55\textwidth}
\begin{block}{六种Mamba学生}
\centering\footnotesize
\begin{tabular}{llcc}
\toprule
模型 & 编码器 & 参数量 & 蒸馏碰撞 \\
\midrule
A & SS2D扫描 & 0.97M & 3 \\
B & MambaVision & 2.61M & 2 \\
\HL{B+} & MambaVision & 2.55M & \HL{1} \\
C & CNN & 2.41M & 3 \\
D & CNN类 & 2.60M & 2 \\
\HL{E} & 轻量CNN & 2.19M & \HL{1} \\
\midrule
教师 & ViT & 3.56M & 2 \\
\bottomrule
\end{tabular}
\end{block}

\column{0.45\textwidth}
\begin{block}{跨架构蒸馏}
\centering\footnotesize
\begin{tikzpicture}[node distance=0.3cm,font=\tiny]
\node[draw,rounded corners,fill=black!10,minimum width=4cm] (T)
{\textbf{ViT+LSTM} $\to$ $\mathcal{L}=\alpha L_f+\beta L_d+\gamma L_g$};
\node[draw,rounded corners,fill=emphasis!10,minimum width=4cm,below=of T] (S)
{\textbf{Mamba学生} (6种)};
\draw[->,thick] (T) -- node[right]{\tiny 三分量损失} (S);
\end{tikzpicture}
\end{block}

\begin{block}{对照基线}
\centering\footnotesize
G\_basic (CNN+MLP, 0.49M)\\
G\_lstm (CNN+LSTM, 0.80M)
\end{block}
\end{columns}
\end{frame}

% ═══════ 4. FIRST SURPRISE ═══════
\begin{frame}{第一个意外：推理速度的反直觉结果}
\begin{columns}[T]
\column{0.5\textwidth}
\begin{block}{预期 vs 现实}
\begin{itemize}
\item \textbf{初始假设}：参数量小 = 推理快
\item \textbf{实际结果}：\HL{A (0.97M)} 推理 \con{24.3ms}
\item \textbf{根因}：SS2D编码器四方向扫描，每次独立执行
\end{itemize}
\end{block}

\begin{block}{关键认知修正}
\footnotesize
\textbf{参数量 $\neq$ 推理速度}\\
编码器架构对延迟的影响远大于时序头
\end{block}

\column{0.5\textwidth}
\begin{table}
\centering\footnotesize
\begin{tabular}{lcc}
\toprule
模型 & 参数量 & 推理延迟 \\
\midrule
\HL{G\_basic} & 0.49M & \HL{0.74ms} \\
G\_lstm & 0.80M & 1.00ms \\
\HL{E} & 2.19M & 7.1ms \\
教师 & 3.56M & 9.0ms \\
B+ & 2.55M & 9.8ms \\
\con{A} & \con{0.97M} & \con{24.3ms} \\
\bottomrule
\end{tabular}
\end{table}
\footnotesize\centering G\_basic快9.6倍，但碰撞多1.8次
\end{columns}
\end{frame}

% ═══════ 5. THE CORE DISCOVERY ═══════
\begin{frame}{核心发现：蒸馏有效，但有硬性条件}
\begin{columns}[T]
\column{0.48\textwidth}
\begin{block}{蒸馏效果 (球体, 5m/s)}
\centering\footnotesize
\begin{tabular}{lcc}
\toprule
模型 & BC & 蒸馏 \\
\midrule
\HL{B+} & 3 & \HL{1}$\downarrow$ \\
\HL{E} & 3 & \HL{1}$\downarrow$ \\
B & DNF & 2$\uparrow$ \\
D & 2 & 2 \\
C & 3 & 3 \\
A & 3 & 3 \\
\bottomrule
\end{tabular}
\vspace{0.1cm}
B+和E蒸馏后降到1次碰撞，\pos{超越教师的2次}
\end{block}

\column{0.52\textwidth}
\begin{alertblock}{为什么A和C蒸馏无效？—— E-SSM对照}
\footnotesize
\begin{itemize}
\item A (SS2D)：扫描展平2D$\to$1D。C (CNN)：卷积保留2D结构
\item \textbf{设计E-SSM}：将E的CNN编码器换成SSM，其余不变
\end{itemize}
\vspace{0.2cm}
\centering
\begin{tabular}{lcc}
\toprule
& $\mathcal{L}_\text{feat}$ & $\mathcal{L}_\text{GT}$ \\
\midrule
E (CNN) & 0.85 & 0.02 \\
\con{E-SSM (SSM)} & \con{10.58} & \con{0.83} \\
\bottomrule
\end{tabular}
\vspace{0.1cm}
\textbf{\con{蒸馏完全失败}}
\end{alertblock}
\end{columns}
\end{frame}

% ═══════ 6. THE "AHA" ═══════
\begin{frame}{"Aha moment"：空间结构是蒸馏的前提}
\begin{columns}[T]
\column{0.5\textwidth}
\begin{block}{为什么CNN有效，SSM失败}
\centering
\begin{tikzpicture}[node distance=0.4cm,font=\footnotesize]
\node[draw,rounded corners,fill=positive!15,text width=4cm,align=center] (C)
{\textbf{CNN编码器}\\[2pt]卷积保留2D空间排列\\$\to$ 与ViT教师特征空间兼容\\$\to$ MSE对齐有效};
\node[draw,rounded corners,fill=negative!15,text width=4cm,align=center,below=of C] (S)
{\textbf{SSM编码器}\\[2pt]选择性扫描展平空间维度\\$\to$ 逐点对应关系被破坏\\$\to$ MSE对齐完全失败};
\end{tikzpicture}
\end{block}

\column{0.5\textwidth}
\begin{block}{上升为一般性原则}
\footnotesize
\textbf{跨架构蒸馏的第一前提：}\\
\textbf{\HL{学生编码器必须保留空间结构}}\\[4pt]
这不仅适用于ViT$\to$Mamba场景。\\
对未来所有跨架构蒸馏的编码器选型\\
都具有直接的指导意义。
\end{block}

\begin{alertblock}{这个发现的特殊性}
\footnotesize
我们没有满足于「蒸馏对某些模型有效」，\\
而是追问了\textbf{为什么}——设计了E-SSM对照实验，\\
找到了因果机制。这才是科学发现。
\end{alertblock}
\end{columns}
\end{frame}

% ═══════ 7. BEYOND EXPECTATION ═══════
\begin{frame}{第二个意外：蒸馏竟让学生超越教师}
\begin{columns}[T]
\column{0.48\textwidth}
\begin{block}{速度鲁棒性 (7m/s)}
\centering\footnotesize
\begin{tabular}{lcc}
\toprule
模型 & 5m/s & 7m/s \\
\midrule
\HL{E 蒸馏} & 1 & \HL{1} \\
\HL{B+ 蒸馏} & 1 & \HL{1} \\
教师 ViT & 2 & \con{5} ($\uparrow$150\%) \\
\bottomrule
\end{tabular}
\end{block}
\vspace{0.2cm}
\footnotesize
蒸馏模型在高速下\textbf{零退化}，\\
教师反而退化150\%。\\
\pos{蒸馏增强了速度鲁棒性，}\\
\pos{超越了教师自身的适应范围。}
\end{block}

\column{0.5\textwidth}
\begin{block}{控制质量与效率权衡}
\centering\footnotesize
\begin{tabular}{lccc}
\toprule
& MAE & Jerk & MAE\_y \\
\midrule
\HL{E} & \HL{0.220} & \HL{0.023} & \HL{0.060} \\
教师 & 0.346 & 0.040 & 0.307 \\
\bottomrule
\end{tabular}
\vspace{0.15cm}
\textbf{但SSM时序头贡献有多大？}\\
G\_basic (CNN+MLP, 0.49M, 0.74ms)\\
仅比E多1.8碰撞，参数少78\%，快9.6倍\\
$\to$ SSM时序头在T=1下边际贡献有限
\end{block}
\end{columns}
\end{frame}

% ═══════ 8. EXPERIMENTAL DESIGN ═══════
\begin{frame}{实验设计：从自动化到系统验证}
\begin{columns}[T]
\column{0.48\textwidth}
\begin{block}{三阶段递进}
\begin{enumerate}\footnotesize
\item \textbf{基线建立}：6种架构 BC 100轮\\
$\to$ 理解各架构的裸性能
\item \textbf{蒸馏引入}：6种架构 蒸馏 50轮\\
$\to$ 测试跨架构知识迁移
\item \textbf{消融验证}：5个维度 20+组\\
$\to$ 分离各因素的影响
\end{enumerate}
\end{block}

\begin{block}{人机协作模式}
\footnotesize
AI：代码实现·图表绘制·文本撰写\\
人类：架构决策·实验设计·结果解释\\
关键：\HL{人类做决策，AI做执行}
\end{block}

\column{0.48\textwidth}
\begin{block}{双Agent流水线}
\centering
\begin{tikzpicture}[node distance=0.3cm,font=\footnotesize]
\node[draw,rounded corners,fill=positive!15,text width=3.5cm,align=center] (T)
{\textbf{训练Agent}\\[2pt]\tiny AutoDL·RTX 5090};
\node[draw,rounded corners,fill=emphasis!15,text width=3.5cm,align=center,below=of T] (S)
{\textbf{仿真Agent}\\[2pt]\tiny WSL2·Flightmare};
\draw[->,thick] (T.east) to[out=0,in=0] node[right]{\tiny GitHub} (S.east);
\draw[->,thick] (S.west) to[out=180,in=180] node[left]{\tiny 结果} (T.west);
\end{tikzpicture}
\end{block}

\begin{block}{实验规模}
\centering\footnotesize
48组主实验 + 20+组消融\\
60+轮训练，100+次仿真\\
球体和树木两种环境，5/7m/s两种速度
\end{block}
\end{columns}
\end{frame}

% ═══════ 9. WHAT WENT WRONG ═══════
\begin{frame}{诚实面对：我们做不到的}
\begin{columns}[T]
\column{0.5\textwidth}
\begin{block}{方法论局限}
\begin{enumerate}\footnotesize
\item \textbf{仅在仿真中验证}\\
Flightmare合成深度图 $\neq$ RealSense真实噪声\\
未部署到真实四旋翼
\item \textbf{种子统计不全}\\
仅G\_basic有4种子统计\\
其余模型结果来自单次训练
\item \textbf{评估维度单一}\\
碰撞次数是主要指标\\
缺少能耗、偏航稳定性等
\end{enumerate}
\end{block}

\column{0.5\textwidth}
\begin{block}{范式局限}
\begin{enumerate}\footnotesize
\item \textbf{BC的固有限制}\\
协变量偏移未根治\\
蒸馏缓解但未消除
\item \textbf{蒸馏的泛化边界}\\
树木环境 7m/s：\\
教师0碰撞，E蒸馏2碰撞\\
$\to$ 学生容量限制
\item \textbf{SSM时序头的疑问}\\
G\_basic (纯CNN+MLP)\\
仅比E多1.8碰撞但快9.6倍\\
$\to$ \textbf{SSM时序头到底贡献了多少？}
\end{enumerate}
\end{block}
\end{columns}
\end{frame}

% ═══════ 10. DESIGN PRINCIPLES + TAKEAWAY ═══════
\begin{frame}{我们学到了什么}
\begin{columns}[T]
\column{0.5\textwidth}
\begin{block}{设计因素优先级}
\centering\footnotesize
\begin{tabular}{lcc}
\toprule
因素 & 碰撞范围 & 影响力 \\
\midrule
\HL{编码器类型} & 1--5 & \HL{★★★★★} \\
知识蒸馏 & 1--DNF & ★★★★☆ \\
数据增强 & 1--5 & ★★★☆☆ \\
损失权重 & 1--4 & ★★☆☆☆ \\
多步预测 & 1--5 & ★★☆☆☆ \\
\bottomrule
\end{tabular}
\end{block}
\vspace{0.2cm}
\begin{itemize}\footnotesize
\item T=1单步始终最优 — 当前帧精确感知 $>$ 长程记忆
\item 数据增强 \pos{CNN有益} 但 \con{SSM有害}
\end{itemize}

\column{0.5\textwidth}
\begin{block}{可操作的建议}
\begin{enumerate}\footnotesize
\item \textbf{选编码器优先于选时序头}\\
编码器质量主导性能
\item \textbf{蒸馏前先确认空间结构兼容性}\\
CNN/混合编码器 $\checkmark$\\
纯SSM编码器 $\times$
\item \textbf{部署配置推荐}\\
E蒸馏：2.19M·7.1ms·1碰撞·速度鲁棒\\
需泛化性时选用B+ (MambaVision)
\end{enumerate}
\end{block}
\end{columns}
\end{frame}

% ═══════ 11. SUMMARY ═══════
\begin{frame}{总结：从问题到答案的路程}
\begin{enumerate}
\item \textbf{起点}：Mamba能替代Transformer吗？——自动化探索效果不好
\item \textbf{转折}：AI建议蒸馏，转向「怎么让Mamba更好」
\item \textbf{核心发现}：\HL{跨架构蒸馏有效，但编码器必须保留空间结构}\\
\footnotesize ——通过E-SSM对照实验确立了因果机制
\item \textbf{意外收获}：\HL{蒸馏使学生超越教师}\\
\footnotesize 高速下教师退化150\%，蒸馏模型零退化
\item \textbf{反思}：\HL{SSM时序头在单帧决策中贡献有限}\\
\footnotesize G\_basic (CNN+MLP) 仅比E多1.8碰撞，但快9.6倍
\end{enumerate}
\vspace{0.3cm}
\begin{center}
\large\textbf{六种架构·四十篇文献·六十轮训练·一百次仿真}\\
\textbf{两个意外发现·一条设计原则·一个可操作指南}
\end{center}
\end{frame}

% ═══════ REFERENCES ═══════
\begin{frame}{参考文献}
\footnotesize
\begin{thebibliography}{9}
\bibitem{gu2023mamba} Gu \& Dao. Mamba: Linear-time sequence modeling with selective state spaces. \textit{arXiv}, 2023.
\bibitem{dao2024mamba2} Dao \& Gu. Transformers are SSMs. \textit{arXiv}, 2024.
\bibitem{bhattacharya2024vitfly} Bhattacharya et al. ViTs for quadrotor obstacle avoidance. \textit{ICRA}, 2025.
\bibitem{liu2024vmamba} Liu et al. VMamba: Visual state space model. \textit{arXiv}, 2024.
\bibitem{hatamizadeh2024mambavision} Hatamizadeh \& Kautz. MambaVision. \textit{arXiv}, 2024.
\bibitem{hinton2015distilling} Hinton et al. Distilling knowledge in a neural network. \textit{NeurIPS Workshop}, 2015.
\bibitem{bick2024mohawk} Bick et al. MOHAWK: Transformers to SSMs. \textit{NeurIPS}, 2024.
\bibitem{shah2017flightmare} Shah et al. Flightmare: A flexible quadrotor simulator. \textit{CoRL}, 2017.
\bibitem{loquercio2018dronet} Loquercio et al. DroNet: Learning to fly by driving. \textit{IEEE RA-L}, 2018.
\end{thebibliography}
\end{frame}

\begin{frame}[plain]
\begin{center}
\Huge 谢谢！\\[1cm]
\Large 欢迎提问\\[1.5cm]
\normalsize GitHub: \texttt{github.com/Liber1917/vitfly}
\end{center}
\end{frame}

\end{document}
